\documentclass[11pt]{article}

\usepackage[margin=1in]{geometry}
\usepackage[T1]{fontenc}
\usepackage[utf8]{inputenc}
\usepackage{lmodern}
\usepackage{microtype}
\usepackage{amsmath, amssymb}
\usepackage{booktabs}
\usepackage{xcolor}
\usepackage{hyperref}
\usepackage{enumitem}
\usepackage{graphicx}
\usepackage{float}

\hypersetup{
  colorlinks=true,
  linkcolor=blue!50!black,
  urlcolor=blue!50!black,
  citecolor=blue!50!black
}

\title{Context versus Chunking:\\A Hidden Signed-Velocity Control Toy}
\author{Andy Park\\\href{mailto:andypark.purdue@gmail.com}{andypark.purdue@gmail.com}}
\date{August 18, 2026}

\begin{document}
\maketitle

\begin{abstract}
This note tests a narrow control mechanism: longer temporal context can denoise inference of a persistent hidden signed-velocity regime, while action-chunk horizon controls how often the policy can refresh after surprises. The final sweep contains 25 context/horizon cells and 120 paired trials per cell. Persistent-mode accuracy rises from 57.8\% for one displacement to 100.0\% for 16 displacements in the fixed persistent-regime check. With context held at C16, the mean paired H16-minus-H1 restricted event-to-recovery gap is 2.32 steps, with a deterministic bootstrap 95\% interval of [1.29, 3.37]. This is a synthetic analytical-controller experiment, not a trained VLA or a real-robot claim.
\end{abstract}

\section{Motivation}

Observation history and action chunking solve different problems. History can average observation noise and reveal motion that is ambiguous in a single displacement. Chunking amortizes planning calls and can smooth actions, but a committed open-loop chunk cannot react until the next planning boundary. This experiment places both mechanisms in one paired toy without assuming that one configuration must dominate every metric.

\section{References and Scope}

The experiment follows standard least-squares line fitting for the temporal estimator and standard feedback-control ideas for the analytical tracker \cite{nist-ls,feedback}. It is not a reproduction of either reference, a VLA paper, or a learned action-chunking system. All numerical results below are local measurements generated by the repository script.

\section{Toy Setup}

A one-dimensional target has hidden mode $m_t\in\{-1,+1\}$. During a persistent regime its velocity is
\[
  v_t^\star = m_t v_0 + \epsilon_t,
\]
where $\epsilon_t$ is correlated zero-mean process noise. The mode flips with a fixed per-step probability. The robot receives only noisy target position $y_t=p_t^\star+\eta_t$; hidden mode, target velocity, and disturbances are retained only for analysis.

At each planning call, C$k$ fits a line to the latest $k+1$ positions, i.e. the latest $k$ temporal displacements. C1 therefore uses $y_t-y_{t-1}$ and is a meaningful but noisy baseline. The slope $\hat v_t$ is clipped to $[-1.35,1.35]$ to bound early partial-window estimates. A nominal PD tracker then predicts target position using $\hat v_t$ and emits an H-step action chunk. Those H actions execute open loop unless the episode ends.

The sweep crosses $C,H\in\{1,2,4,8,16\}$. For each trial seed, every cell receives exactly the same hidden-mode sequence, correlated target-velocity noise, observation noise, and robot impulses. Twelve separate closed-loop oracle demonstrations are also generated. Their oracle actions use true target state, but the evaluated controller is analytical and is \textbf{not trained from the demonstration CSV}.

\section{Metrics}

Let $e_t=|p_t^\star-p_t|$ denote error after the mode-dependent target transition and robot impulse at step $t$.
\begin{itemize}[leftmargin=1.5em]
  \item \textbf{Success:} at least 73\% of the final 45 errors satisfy $e_t<0.68$.
  \item \textbf{Tracking error:} episode mean $T^{-1}\sum_t e_t$.
  \item \textbf{Persistent-mode accuracy:} $\operatorname{sign}(\hat v_t)=m_t$ at planning calls with a complete C-transition window and at least 16 steps since the most recent switch. The common persistence guard separates denoising from unavoidable post-switch window lag.
  \item \textbf{Jerk proxy:} $\frac{1}{T-1}\sum_{t=1}^{T-1}(u_t-u_{t-1})^2/\Delta t^2$. This is an action-difference proxy, not physical jerk.
  \item \textbf{Planning calls:} number of analytical policy invocations per episode; it is not measured wall-clock latency.
\end{itemize}

\subsection{Event recovery definition}

For an event at step $s$, the event is applied before $e_s$ is recorded. The baseline $b_s$ is the median of the preceding 6 errors. The response window is $[s,n)$, where $n$ is the next event or episode end. Within its first 10 samples, let $q$ be the post-event peak at step $p$. The event is \emph{impactful} only when $q-b_s\ge 0.10$. Threshold search starts strictly after $p$, where the recovery condition is necessarily false, and ends at the first later pair of consecutive errors below
\[
  b_s + 0.25(q-b_s).
\]
The primary delay is measured from the event step $s$, not from the peak; peak-to-recovery delay is retained only as a diagnostic. If no such pair appears before $n$, the event is unrecovered and contributes the censoring delay $n-s$ to the restricted mean. Event coverage is impactful divided by evaluable scheduled events; unrecovered counts are reported explicitly. Events without enough pre-event history or response room are not evaluable, and events without a demonstrated rise are excluded rather than counted as zero-delay recoveries.

\section{Results}

The committed artifacts were generated with:
\begin{verbatim}
cd /home/andypark/Projects/repos/vla-ideas
/home/andypark/Projects/repos/dhb-xr/.pixi/envs/default/bin/python \
  context_chunk_tradeoff/run_context_chunk_tradeoff.py \
  --trials 120 --seed 17
\end{verbatim}

\begin{table}[H]
\centering
\resizebox{\textwidth}{!}{%
\begin{tabular}{lrrrrrrrr}
\toprule
Controller & Success & Error & Mode acc. & Recovery & Coverage & Unrec. & Jerk & Calls \\
\midrule
C1 / H1 & 44.2\% & 0.588 & 57.2\% & 18.14 & 45.5\% & 188/307 & 795.7 & 180 \\
C16 / H1 & 50.0\% & 0.479 & 100.0\% & 14.01 & 69.3\% & 237/467 & 21.3 & 180 \\
C16 / H16 & 30.0\% & 1.028 & 100.0\% & 15.48 & 67.8\% & 255/457 & 32.1 & 12 \\
\bottomrule
\end{tabular}}
\caption{Selected baselines from 120 paired trials per cell (seed 17). Recovery is the pooled restricted mean over impactful events. Coverage is impactful/evaluable; Unrec. is unrecovered/impactful.}
\label{tab:baselines}
\end{table}

The deterministic persistent-regime check gives C1 mode accuracy 57.8\% and C16 accuracy 100.0\%. Hand-computed traces verify post-peak threshold search, no-rise exclusion, and next-event censoring. The small fixed surprise schedule gives 13.64 event-to-recovery steps for H1 and 13.79 for H16, but is descriptive rather than a pass criterion. Across the full paired C16 sweep, the per-seed H16-minus-H1 gap is 2.32 steps (deterministic bootstrap 95\% interval [1.29, 3.37]).

\begin{figure}[H]
  \centering
  \includegraphics[width=\textwidth]{../outputs/tradeoff_heatmaps.png}
  \caption{The full 25-cell sweep. Mode-inference and jerk evidence are shown directly alongside task, recovery, and planning metrics. Annotation color adapts to the cell luminance.}
  \label{fig:heatmaps}
\end{figure}

\begin{figure}[H]
  \centering
  \includegraphics[width=0.88\textwidth]{../outputs/pareto_tradeoff.png}
  \caption{Reactivity--smoothness trade-off. Longer chunks reduce planning calls and reduce the action-difference jerk proxy on average across contexts. With context held at C16, H1 has lower restricted event-to-recovery delay than H16.}
  \label{fig:tradeoff}
\end{figure}

\begin{figure}[H]
  \centering
  \includegraphics[width=\textwidth]{../outputs/paired_rollout.png}
  \caption{One paired exogenous rollout. Estimates are bounded during warm-up. Gray lines mark mode switches or impulses; the H16 controller can retain a stale plan between boundaries.}
  \label{fig:rollout}
\end{figure}

\section{Interpretation}

In the intended persistent regime, longer context materially improves inference because it averages noisy position displacements. The same history can lag immediately after a switch because old-regime samples remain in the window, so context is not claimed to help monotonically at every instant. Longer action chunks retain the expected planning-call advantage and reduce the action-difference jerk proxy on average across context settings, but smoothness is not monotonic at fixed C16. With context held at C16, H1 has lower restricted event-to-recovery delay than H16 because it can replace stale plans sooner; that reactivity claim is not generalized to every context setting. These are metric trade-offs, not an encoded automatic winner.

\clearpage
\section{Limitations and Follow-ups}

The tracker is one-dimensional, its policy is analytical, and train/test distribution shift is absent because there is no training stage. Position noise is Gaussian; hidden modes are binary; disturbances are velocity impulses; and the controller has no images, language, contact, latency, or hardware. Recovery is conditional on a measurable error response, so coverage and unrecovered accounting are essential to interpretation. The response threshold, persistence guard, and success tolerance are local definitions rather than community benchmarks.

Useful follow-ups are to pre-register the recovery constants on a calibration set, evaluate more switch frequencies and non-Gaussian noise, separate sensing from planning latency, replace position observations with learned visual features, and repeat the paired design with a trained chunking policy.

\begin{thebibliography}{9}
\bibitem{nist-ls} NIST/SEMATECH. \emph{Linear Least Squares Regression}. \url{https://www.itl.nist.gov/div898/handbook/pmd/section1/pmd141.htm}
\bibitem{feedback} K. J. \AA str\"om and R. M. Murray. \emph{Feedback Systems: An Introduction for Scientists and Engineers}. Princeton University Press, 2008. \url{https://fbswiki.org/wiki/index.php/Main_Page}
\end{thebibliography}

\end{document}
